\subsection{جابه‌جایی دامنه‌ی بازنمایی‌ها و نشانگرها}

پرسش این زیربخش توصیفی است. هر منبع کنارگذاشته در فضای بازنمایی‌های منجمد چه اندازه از منابع آموزشی تجمیع‌شده فاصله دارد، و این فاصله میان بلوک تصویری، بلوک عروقی و پنج نشانگر اسکالر چگونه پخش می‌شود؟ معیار، تفاوت میانگین استانداردشده‌ی هر ویژگی (\lr{SMD}\enfoot{تفاوت میانگین استانداردشده‌ی یک ویژگی میان دو جمعیت؛ معادل انگلیسی: \lr{Standardized Mean Difference}}) است، یعنی میانگین منبع کنارگذاشته منهای میانگین منابع آموزش، تقسیم بر انحراف معیار منابع آموزش. این آماره هیچ برچسب هدفی مصرف نمی‌کند، هیچ آزمون استنباطی در آن اجرا نشده و هیچ مدلی با آن بازتنظیم نشد.

\begin{table}[htbp]
\centering
\renewcommand{\arraystretch}{1.35}
\small
\caption{خلاصه‌ی جابه‌جایی دامنه‌ی نمایش‌های تثبیت‌شده؛ منبع کنارگذاشته در برابر منابع آموزشی تجمیع‌شده}
\label{tab:domainshift}
\setlength{\tabcolsep}{3pt}
\begin{tabular}{lrrrrr}
\hline
منبع کنارگذاشته & نمایش & میانه‌ی قدر مطلق \lr{SMD} & صدک 90 & صدک 95 & فاصله‌ی مرکزوار استانداردشده \\
\hline
\lr{Plus} & تصویری & 0.5406 & 1.8314 & 2.5360 & 56.511 \\
\hline
\lr{Plus} & عروقی & 1.2291 & 2.3345 & 2.6296 & 63.404 \\
\hline
فارفوم & تصویری & 0.3793 & 0.8431 & 1.1086 & 26.038 \\
\hline
فارفوم & عروقی & 0.5565 & 0.7307 & 0.7579 & 22.953 \\
\hline
فارابی & تصویری & 0.4966 & 1.1258 & 1.4847 & 35.253 \\
\hline
فارابی & عروقی & 0.6972 & 0.8829 & 0.9280 & 28.324 \\
\hline
\multicolumn{6}{p{0.96\textwidth}}{\footnotesize \lr{SMD} هر ویژگی برابر تفاضل میانگین منبع کنارگذاشته و میانگین منابع آموزشی تقسیم بر انحراف معیار منابع آموزشی است. هیچ برچسب هدفی در این آماره‌ها استفاده نشده و هیچ مدلی با آن‌ها بازتنظیم نشد. نمایش عروقی در هر سه فولد جابه‌جاشده‌تر از نمایش تصویری است، ولی همان نمایش عروقی بیشترین سهم را در بهبود رتبه‌بندی دارد.} \\
\hline
\end{tabular}
\end{table}

جدول میانه و صدک‌های قدر مطلق \lr{SMD} را برای دو بلوک نشان می‌دهد. بر پایه‌ی این میانه، بلوک عروقی در هر سه فولد جابه‌جاشده‌تر از بلوک تصویری است، در حالی که همین بلوک بیشترین سهم را در بهبود رتبه‌بندی دارد؛ افزودن آن به بازنمایی تصویری در دو فولد سه‌کلاسه‌ی منبع‌کنارگذاشته تفاضل مثبت و معنادار داد. پس حساسیت به منبع به‌خودی‌خود آن چیزی نیست که یک بازنمایی را از سودمندی بازمی‌دارد.

دو آماره‌ی جدول یک روایت واحد نمی‌سازند. فاصله‌ی مرکزوار یک جمع سراسری روی همه‌ی ابعاد است و میانه‌ی قدر مطلق \lr{SMD} توزیع تک‌ویژگی را می‌سنجد، بنابراین این دو عدد لزوماً هم‌جهت نیستند و می‌توانند در یک فولد دو مسیر متفاوت بروند.

\begin{figure}[htbp]
  \centering
  \imgbox[1.0\textwidth]{figures/report/main/DS1a_shift_rgb.pdf}
  \caption{جابه‌جایی دامنه‌ی نمایش تصویری میان منبع کنارگذاشته و منابع آموزشی.}
  
\end{figure}

\begin{figure}[htbp]
  \centering
  \imgbox[1.0\textwidth]{figures/report/main/DS1b_shift_vessel.pdf}
  \caption{جابه‌جایی دامنه‌ی نمایش عروقی میان منبع کنارگذاشته و منابع آموزشی.}
  \label{fig:U25}
\end{figure}

\begin{figure}[htbp]
  \centering
  \imgbox[1.0\textwidth]{figures/report/main/DS2_centroid_shift.pdf}
  \caption{فاصله‌ی مرکزوار استانداردشده میان منبع کنارگذاشته و منابع آموزشی در دو نمایش.}
  \label{fig:U26}
\end{figure}

\begin{figure}[htbp]
  \centering
  \imgbox[1.0\textwidth]{figures/report/main/DS3_biomarker_smd_heatmap.pdf}
  \caption{تفاوت میانگین استانداردشده‌ی هر نشانگر میان منبع کنارگذاشته و منابع آموزشی.}
  \label{fig:U27}
\end{figure}




پنج نشانگر اسکالر هم از این جابه‌جایی بی‌بهره نیستند و اندازه‌ی جابه‌جایی آن‌ها تا حدود یک انحراف معیار کامل داده‌ی آموزش می‌رسد. علامت‌ها میان منابع برمی‌گردند؛ مثبت در \lr{Plus} و منفی در دو منبع دیگر. یعنی یک نشانگر با تعریف ثابت در سه مرکز نه عدد یکسان و نه جهت یکسان به دست نمی‌دهد (شکل شکل مربوط).



این اعداد نشان می‌دهند بازنمایی‌ها و نشانگرها هر دو به منبع تصویربرداری حساس‌اند و بلوکی که بیشترین سهم را دارد، جابه‌جاشده‌ترین بلوک است. از این هم‌زمانی هیچ ادعای علّی ساخته نمی‌شود و هیچ داوری درباره‌ی کیفیت یا بی‌کیفیتی یک بازنمایی از آن درنمی‌آید. آنچه با اطمینان گفته می‌شود دو چیز است. حساسیت به منبع به‌تنهایی یک بازنمایی را از سودمندی بازنمی‌دارد، و پنج نشانگر اسکالر با این اندازه از جابه‌جایی برای انتقال بی‌تغییر به یک مرکز تازه مناسب نیستند و مقیاس ثابت آن‌ها میان مراکز جابه‌جا می‌شود.
